\section{The integral complementary basis and its exact count}
\label{sec:complementary-basis}

\status{This supplies the integral basis, complementary-family count and Cartan multiplication used in CL1--CL5. The proofs work over every integral domain, including residue fields; they use polynomial identities, not evaluations at the points of a finite field.}

\paragraph{CB1. A leading-term calculation.}
Let $R$ be a domain and $X$ an $r\times n$ matrix of indeterminates.
Write $\nabla_\nu(R)$ for polynomials of column multidegree $\nu$ invariant under every generic operation $C_j\mapsto C_j+tC_l$, $l<j$.
Order variables column first, then row, with earlier variables larger.
If $X^e$ leads a nonzero member, then
\[
 \sum_a e_{a,t}=\nu_t,\qquad
 \sum_{b\le a}e_{b,j}\le\sum_{b<a}e_{b,j-1}. \tag{CB.1}
\]
Here is the substantive step. For two columns $x,y$, the more general bound
$\deg_t f(x,y+tx)\le s$ implies
$\sum_{i\le a}b_i\le s+\sum_{i<a}a_i$ for their leading exponents.
First retain the homogeneous component fixing each pair's total degree and every other variable's exponent; this preserves the leading monomial and the bound.
At the first row write
\[
 f=\sum_{u=0}^{K}x^uy^{e+K-u}f_u,\qquad f_K\ne0.
\]
After substitution this is $(y+tx)^eQ$, with
$\deg_tQ\le s-e$ and $e\le s$.
Setting $x=1,y=-t$ in $Q$ leaves $f_K(x',y'+tx')$; since $\deg_yQ\le K$, its $t$-degree is at most $s-e+K$.
Induction on rows proves the assertion. Generic invariance gives $s=0$ and hence (CB.1).

\paragraph{CB2. A monic integral basis.}
For an array satisfying (CB.1), make tableau row $t$ from $e_{a,t}$ copies of $a$, in increasing order.
The inequalities say exactly that entries increase strictly down columns.
For each tableau column take the minor using its entries as matrix rows and the first matrix columns; multiply these minors to obtain $b_e$.
The diagonal term leads each minor with coefficient one, so $b_e$ has leading monomial $X^e$.
Every $b_e$ belongs to $\nabla_\nu(R)$, and distinct $e$ give independence.
For any $f$, subtract its leading coefficient times the corresponding $b_e$.
The leading monomial decreases; the fixed multidegree admits finitely many monomials, so this terminates.
Thus the $b_e$ are an integral basis, homogeneous for row weights.
The same integer polynomials give bases after passage to any domain, including reduction modulo a prime.
Products in disjoint blocks likewise give a tensor basis.

\paragraph{CB3. Counting without a dimension theorem.}
For a decreasing integer tuple $\lambda$ of length $k$, put
\[
 D_k(\lambda)=\prod_{i<j}\frac{\lambda_i-\lambda_j+j-i}{j-i},\qquad D_0=D_1=1.
\]
Tableaux of shape $\nu=\lambda-\lambda_k$ with entries at most $k$ have this count.
Indeed, deleting the entries $k$ leaves a shape $\mu$ satisfying
$\nu_i\ge\mu_i\ge\nu_{i+1}$.
Induct on $k$, using $x_i=\nu_i+k-i$ and $y_i=\mu_i+k-1-i$ and the identity
\[
 \sum_{x_{i+1}\le y_i<x_i}\prod_{i<j}(y_i-y_j)
 =\frac{1}{(k-1)!}\prod_{i<j}(x_i-x_j). \tag{CB.2}
\]
For completeness, write the left Vandermonde as a determinant and sum each row.
Choose $Q_j$ with $Q_j(t+1)-Q_j(t)=t^{k-1-j}$, of leading coefficient $1/(k-j)$.
The resulting determinant is $\det[Q_j(x_i)-Q_j(x_{i+1})]$.
Append the constant column and undo successive row differences: the determinant has degrees $k-1,\ldots,0$ and leading-coefficient product $1/(k-1)!$, proving (CB.2).

\paragraph{CB4--CB5. The actual complement.}
For a selected weight $\lambda^0$, put
\[
 \ell=(\lambda^0_2,\lambda^0_2,\lambda^0_3,\ldots,\lambda^0_k).
\]
Use the effective weight $\ell$ at even positions and $-\operatorname{rev}\ell$ at odd positions, and subtract its last entry to obtain $\nu$.
At an even position impose $e_{1,1}=\nu_1$: deleting the first row and column gives shape $(\nu_2,\ldots,\nu_k)$.
At an odd position impose $e_{1,t}=0$: since $\nu_k=0$, delete the first row and last column.
The inverse operations insert the indicated zeros and, in the even case, $\nu_1$.
Constant shifts and negative reversal preserve $D$, so both counts are $D_{k-1}(\lambda^0_2,\ldots,\lambda^0_k)$.
Consequently
\[
 n_2=\prod_{s\in T}D_{k-1}(\lambda^0_{s,2},\ldots,\lambda^0_{s,k})
       \prod_{s\notin T}D_k(\lambda^0_s). \tag{CB.3}
\]
The tensor basis also gives $\dim V=\prod_sD_k(\lambda^0_s)$, the source's \texttt{FewBadOrbits} interface.
If $d_s=\lambda^0_{s,1}-\lambda^0_{s,2}\ge a>0$ and
$0\le\lambda^0_{s,2}-\lambda^0_{s,j}\le c_0$, then
\[
\begin{aligned}
 \frac{\binom{d_s+k-1}{k-1}D_{k-1}(\lambda^0_{s,2},\ldots,\lambda^0_{s,k})}{D_k(\lambda^0_s)}
 &=\prod_{j=2}^k\frac{d_s+j-1}{d_s+\lambda^0_{s,2}-\lambda^0_{s,j}+j-1}\\
 &\ge1-\frac{(k-1)c_0}{a}.
\end{aligned}
\]
Multiplication yields the source's weaker factor $(1-kc_0/a)^{|T|}$ when $kc_0/a\le1$.

\paragraph{CB6--CB7. Integral multiplication and the source lattices.}
A product of minors of shape $\nu+\mu$ has
$(\nu_j-\nu_{j+1})+(\mu_j-\mu_{j+1})$ factors of height $j$.
Partition them accordingly. CB2 therefore proves
\[
 \nabla_{\nu+\mu}(R)=\operatorname{span}_R
 \{fg:f\in\nabla_\nu(R),\ g\in\nabla_\mu(R)\}. \tag{CB.4}
\]
Over a DVR primitive factors have primitive product: reduce their coefficient polynomials modulo the uniformizer.

For the actual basis matrix $M$, row substitution $\mathcal R(M)f(X)=f(MX)$ is an anti-action. The source lattice is exactly
\[
 L_\lambda(M)=\det(M)^{-\lambda_k}\mathcal R(M^{-1})\nabla_{\lambda-\lambda_k}(\mathcal O).
\]
Changing $M$ to $MU$, $U\in\mathrm{GL}_k(\mathcal O)$, preserves this lattice; hence it agrees with the source's intersection over all bases.
For the midpoint matrices $M_s$ of EL, the actual edge-coordinate isomorphism is
\[
 \mathcal C_\lambda(F)=w^{A_r(\lambda)}
 \left(\prod_s\det(M_s)^{\lambda^{\rm eff}_{s,k}}\right)
 \operatorname{blockAct}(M_s)F.
\]
It takes $R_e(\lambda)$ onto the integral tensor coefficient lattice.
All exponents add with weights, so
$\mathcal C_{\lambda+\mu}(FG)=\mathcal C_\lambda(F)\mathcal C_\mu(G)$.
Thus (CB.4) is the integral edge statement \texttt{OntoAll}, and applying
$\mathcal C_{\rm rest}^{-1}$ to the selected basis products constructs the exact complementary family.
No denominator or characteristic restriction is introduced.
This supplies the basis, count and Cartan inputs retained in CL.
The distinct incoming argument is IN/WT; selected-symbol density and arithmetic windows remain separate.
